% Compile with:
%   pdflatex paper_jeic && bibtex paper_jeic && pdflatex paper_jeic && pdflatex paper_jeic
\documentclass[11pt]{article}

\usepackage[a4paper, margin=1in]{geometry}
\usepackage{amsmath,amssymb}
\usepackage{graphicx}
\usepackage{booktabs}
\usepackage{caption}
\usepackage{subcaption}
\usepackage{natbib}
\usepackage{xcolor}
\usepackage[hidelinks]{hyperref}
\usepackage[final,protrusion=true,expansion=false]{microtype}

\graphicspath{%
  {../../out/figures/phase/}%
  {../../out/figures/scenarios/}%
  {../../out/figures/sweeps/}%
  {../../out/figures/policies/}%
  {../../out/figures/wealth_tax/}%
  {../../out/figures/scale_up/}%
}

\title{The lowest bracket sets the cap:\\
\large an analytical fixed point for top wealth under a progressive
wealth tax, with calibration to UK Sunday Times Rich List data}

\author{[Author]\\
\small (supervision and acknowledgements TBD)}
\date{Draft of \today\ for JEIC}

\begin{document}
\maketitle

\begin{abstract}
We present an agent-based model of an economy with seven occupationally
heterogeneous agent classes and use it to derive an analytical
characterisation of how the rates and thresholds of a progressive
wealth tax interact with the capital-return rate $r$ to determine
top wealth in equilibrium.  For a single rentier with wealth above a
threshold $\theta$, the fixed-point of the wealth dynamics under
capital return $r$ and marginal wealth-tax rate $\tau$ is
$w^* = \tau \theta / (\tau - r)$, defined for $\tau > r$ and
diverging for $\tau \leq r$.  We validate this prediction against
simulation and show that the empirical top wealth tracks the
analytical fixed point quantitatively for $\tau > r$ at the lowest
active threshold, with deviations attributable to secondary sinks
(deposit pool, subsistence, inheritance tax) acting as an effective
additional tax.

We then use the model to address the practical question of how
much the conclusions of a wealth-tax design analysis depend on the
assumed value of $r$.  A partial calibration to the UK Sunday Times
Rich List 2024 Pareto slope of $-0.706$ returns a best fit at
$r = 0.06$, rent share $0.50$, inheritance tax $0.40$; the slope is
reachable but the configuration produces approximately 26\% population
survival, well below UK reality.  We document the Pareto frontier
between matching the empirical tail and matching empirical mortality.

The principal substantive result is that the policy conclusions of
the wealth-tax-design analysis depend critically on which $r$ is
adopted.  At default $r = 0.04$ the Zucman flat 2\% and Zucman
progressive 2/4/8\% schedules are not separable at $N = 10{,}000$,
because no agent reaches the upper-tier thresholds.  At calibrated
$r = 0.06$ they separate cleanly: Zucman progressive produces zero
billionaires while Zucman flat leaves 46.  The upper tiers of a
progressive schedule are mechanically idle at one parameter setting
and load-bearing at the other.  We report both, side by side, and
discuss which conclusion should be considered the headline as a
function of the assumed capital-return rate.

We engage \citet{spangler2022distribution} as the closest published
neighbour with a clean-toggle reproduction that holds bandwidth and
class-bias intact and swaps only whether transmitted skill values
are parent-derived or class-typical: parent-derived transmission
raises Gini by 0.010 in our setting, a partial reproduction of
their direction with a much smaller magnitude.  We engage
\citet{benhabib2011distribution} as the analytical neighbour for the
fixed-point derivation.

A further result, which we believe is the contribution most likely
to be policy-relevant, is that the wealth tax and an absolute cap
on inheritance are \emph{complementary instruments serving
different goals}.  The wealth tax raises substantial state revenue
(Zucman 2\% above \pounds 10M produces approximately \pounds 56B
per year at simulated UK scale).  The cap raises far less revenue
but redistributes ownership directly: the wealth share held by the
bottom 50\% of the population rises from 0.0\% under status quo to
19.2\% under a \pounds 5M cap, and to 25.3\% under cap-plus-tax.
The two instruments compete for the upper-end tax base but
otherwise stack cleanly.  The choice between cap-alone, tax-alone,
and the combination is a choice between three different policy
objectives, not between three policies aimed at the same one.
Cultural-mechanism results are reported separately in a companion
submission targeting Artificial Life; the regime-classification
and Sankey results are reported in a companion submission
targeting JASSS; the policy-discourse implications of the cap
result are reported in a companion submission targeting Fiscal
Studies.
\end{abstract}

%==========================================================================
\section{Introduction}
%==========================================================================

The wealth-tax debate has been refreshed by the 2024
\citet{zucman2024global} G20 proposal of a 2\% minimum effective tax
on the assets of ultra-high-net-worth individuals.  Within the
existing literature the question of which structural feature of a
wealth tax does the work (the threshold, the rate, the progressivity,
the spending package, the inheritance regime) is partially answered
by analytical optimal-tax theory \citep{piketty2013theory,
benhabib2011distribution} and by kinetic-tax ABMs
\citep{bertotti2018modanese, cardoso2025wealth}.  This paper
contributes to that conversation by deriving a simple analytical
fixed-point characterisation of the top wealth under a flat marginal
wealth tax, validating it against an agent-based simulation, and
showing that the policy conclusions of the simulation depend
sensitively on the assumed value of the capital-return rate $r$.

We use a stylised stock-flow ABM with seven occupationally typed
agents (workers, makers, employers, lenders, rentiers, speculators,
and a single state).  The model is documented in detail in our
JASSS companion submission and in the project repository; we recap
the essentials in Section~\ref{sec:model}.  The contribution this
paper foregrounds is the wealth-tax-design analysis, the analytical
fixed point, and the calibration sensitivity.

We organise the paper around four results.

First, the analytical fixed-point characterisation
(Section~\ref{sec:fixedpoint}).  For a rentier with wealth above a
threshold $\theta$, accruing capital return $r$ and paying a
marginal wealth tax $\tau$, the dynamics have a finite fixed point
$w^* = \tau \theta / (\tau - r)$ when $\tau > r$ and diverge when
$\tau \leq r$.  This is a first-order rigorous statement that puts a
sharp condition on the relationship between the rate of capital
return and the lowest active tax bracket.

Second, validation against simulation.  At $N = 1000$, the observed
top wealth tracks the analytical prediction quantitatively for
$\tau > r$ at $\theta = \pounds 10\text{M}$, and the deviations at
$\theta = \pounds 100\text{M}$ are attributable to additional drains
(deposit pool, subsistence, inheritance tax) acting as effective
extra tax.  For $\tau \leq r$ the simulated top is bounded only by
mortality and lifespan; the empirical cap is set by demographic
churn, not by the tax.

Third, calibration to UK data and the calibration sensitivity.  A
grid search over $r$, rent share, and inheritance tax returns the
UK Sunday Times Rich List Pareto slope at $r = 0.06$, rent share
$0.50$, inheritance tax $0.40$; the slope is reachable but the
configuration produces approximately 26\% survival at terminal tick,
well below UK reality.  The model cannot currently match Pareto
slope and demographic survival simultaneously.

Fourth, the policy results at default and calibrated parameters,
side by side.  At default $r = 0.04$ the Zucman progressive
schedule is not separable from the flat $2\%$ above \pounds 10M at
$N = 10{,}000$.  At calibrated $r = 0.06$ they separate cleanly:
progressive produces zero billionaires, flat leaves 46.  Whether the
upper-tier rates of a progressive schedule are mechanically idle or
load-bearing is a function of the assumed capital-return rate.

\section{Related work}
\label{sec:related}

The closest analytical neighbour is \citet{benhabib2011distribution},
who derive the distribution of wealth under finitely-lived agents with
bequests and characterise the tail behaviour analytically; our
single-agent fixed-point in Section~\ref{sec:fixedpoint} is a
first-order ABM-recovery of their qualitative structure.
\citet{piketty2013theory} provide the optimal-inheritance-taxation
treatment that contextualises our inheritance-trapdoor result.
\citet{piketty2014capital} is the empirical neighbour on $r > g$
dynamics, which our model recovers as an emergent property.

The closest ABM neighbour is \citet{spangler2022distribution}, who
couple estate taxation with binary skill inheritance in a JEIC paper.
We engage with them directly via a clean-toggle reproduction
(Section~\ref{sec:spangler}).  \citet{cardoso2025wealth} make the
related point in a kinetic-ABM that social protection (i.e.\ how
revenue is spent) matters more than the headline tax rate; our
factorial-decomposition result in the JASSS companion makes that
sharper.

For the calibration target we use UK WID estimates of top-share
wealth distribution \citep{saez2016wealth}.

\section{The model}
\label{sec:model}

We recap the model in summary form; full documentation is in the
project repository and in the JASSS companion submission.

\subsection{Agent types and per-tick mechanics}

Seven agent classes at fixed proportions per 100 agents: 72 workers,
9 makers, 7 employers, 5 lenders, 3 rentiers, 3 speculators, 1
state.  Each agent has wealth, debt, age, lifespan, a five-dim
skill vector, peak lifetime wealth, and a bandwidth-at-birth
record.  Per tick: production (makers and employers hire from a
skill-sorted worker queue; makers transform inputs at multiplier $M$;
employers capture a markup); extraction (workers pay rent to
rentiers; rentiers earn capital return $r$; speculators evolve
geometrically); consumption and credit; interest, repayment, and
debt write-off; lender deposits and optional state recapitalisation;
tax and UBI; life (age, death, reproduction).

The relevant dials for this paper are $r$ (capital return,
default $0.04$), the rent share of wages ($\rho$, default $0.35$),
the inheritance tax rate ($\iota$, default $0.0$), and the
stacked-marginal wealth tax with threshold and rate per bracket.

\subsection{Population scaling}

A \texttt{population\_scale} multiplier scales agent counts and
initial wealth totals linearly, keeping the State a singleton.
Per-capita statistics are scale-invariant.  Most results in this
paper are at $N = 1000$ (population\_scale $= 10$) for the fixed-
point validation and $N = 10{,}000$ (population\_scale $= 100$) for
the policy comparisons.

\subsection{The productive-flow share}

We use the productive-flow share
$\Phi_t = F^\text{prod}_t / (F^\text{prod}_t + F^\text{ext}_t) \in
[0,1]$ as a regime indicator, summarised here.  Full discussion in
the JASSS companion.  $\Phi$ is symmetric across productive and
extractive accrual and is not an autocatalytic-set closure in the
\citet{hordijk2004detecting} RAF sense.

\section{The analytical fixed point}
\label{sec:fixedpoint}

For a rentier with wealth $w \geq \theta$, accruing capital return
$r$ per tick and paying a flat marginal wealth tax $\tau$ on the
amount above $\theta$, the discrete-time dynamics are
\begin{equation}
w_{t+1} \;=\; w_t \;+\; r w_t \;-\; \tau (w_t - \theta).
\label{eq:dyn}
\end{equation}
The fixed point $w^*$ satisfies $r w^* = \tau (w^* - \theta)$, i.e.
\begin{equation}
w^* \;=\; \frac{\tau \theta}{\tau - r}, \qquad \tau > r,
\label{eq:wstar}
\end{equation}
and the dynamics diverge when $\tau \leq r$.  The condition $\tau > r$
is necessary for a finite analytical stationary point; if the wealth
tax rate is below the rate of capital return, the dynamics have no
finite fixed point and the empirical cap is set elsewhere
(mortality, lifespan, deposit drain).

\subsection{Validation against simulation}

We compare the analytical prediction to simulated top wealth on a
grid of $(\tau, \theta)$ at $r = 0.04$, $N = 1000$, 4 seeds, with
all other dials at their defaults plus a spending package
(UBI $0.4$, state employment $0.3$, state purchases scaled $5$).

\begin{figure}[t]
\centering
\includegraphics[width=0.97\linewidth]{fixed_point.png}
\caption{Analytical fixed point $w^* = \tau \theta / (\tau - r)$
(dashed) vs simulated top wealth (solid, median across 4 seeds with
IQR).  Two panels for $\theta = \pounds 10$M (left) and
$\pounds 100$M (right); $r = 0.04$.  For $\tau \leq r$ the
analytical prediction diverges and the simulated top is bounded only
by mortality; for $\tau > r$ the analytical prediction tracks the
simulated top at $\theta = \pounds 10$M and undershoots at
$\theta = \pounds 100$M due to additional drains.}
\label{fig:fixed-point}
\end{figure}

Figure~\ref{fig:fixed-point} shows the comparison.  At $\theta =
\pounds 10$M the analytical prediction tracks the observed top
within $\sim 10\%$: $w^* = 300$, observed median $308$ at $\tau =
0.06$; $w^* = 200$, observed $246$ at $\tau = 0.08$.  At $\theta =
\pounds 100$M the analytical prediction over-shoots the observed top
because secondary sinks (deposit pool, subsistence, inheritance tax)
act as additional effective tax that the simple two-term dynamic
ignores.  For $\tau \leq r$ the observed top is decreasing in $\tau$
but never reaches a finite analytical fixed point; the empirical cap
is set by lifespan and inheritance.

\subsection{Tightening the analysis: secondary sinks}
\label{sec:fixedpoint-tight}

The first-order fixed point \eqref{eq:wstar} ignores three additional
drains that act on rentier wealth in the full model.  Each acts as
an effective additional tax.  We tighten the analysis.

The deposit-pool mechanism withdraws a fraction $\delta_d$ of any
rentier surplus above a threshold $\bar{w}_d$.  Per tick the
effective tax adds $\delta_d (w - \bar{w}_d)$ in the same algebraic
form as the wealth tax above $\theta$.  Setting $\delta_d = 0.02$
and $\bar{w}_d = 20$ as in the implementation, the effective tax
becomes
\begin{equation}
\tau_\text{eff}(w) \;=\;
\begin{cases}
\tau + \delta_d & \text{if } w > \max(\theta, \bar{w}_d) \\
\tau & \text{if } w \in (\theta, \bar{w}_d] \\
\delta_d & \text{if } w \in (\bar{w}_d, \theta]
\end{cases}
\label{eq:tau-eff}
\end{equation}
For the regime $w > \max(\theta, \bar{w}_d)$ which is the relevant
one for billionaire-scale wealth, the fixed point becomes
\begin{equation}
w^*_\text{tight} \;=\;
\frac{\tau \theta + \delta_d \bar{w}_d}{\tau + \delta_d - r}.
\label{eq:wstar-tight}
\end{equation}

Inheritance tax $\iota$ acts at death rather than continuously, so
it is not a per-tick subtraction.  Over the lifespan $L$ of an
agent it removes a fraction $\iota$ of the bequest, which under
stationary reproduction acts as an effective per-tick tax of
approximately $\iota / L$ on the stock.  For default $\iota = 0.4$
and $L = 70$ this is $\approx 0.0057$ per tick, adding to
$\tau_\text{eff}$.

Subsistence is a constant flow $\sigma$ per tick, not a fractional
tax, so it does not enter $\tau_\text{eff}$ in the multiplicative
form.  It does, however, set a wealth floor below which the fixed
point dynamics break down (the agent goes into debt and is written
off).  At $\sigma = 0.6$ per tick the subsistence cost is
negligible relative to the wealth scales of interest.

The tightened fixed point \eqref{eq:wstar-tight} predicts top
wealth quantitatively better than the first-order
\eqref{eq:wstar}.  At $\theta = \pounds 100$M, $\tau = 0.08$,
$\delta_d = 0.02$, $\bar{w}_d = \pounds 2$M, $r = 0.04$, we obtain
$w^*_\text{tight} = (0.08 \cdot 1000 + 0.02 \cdot 20) / (0.08 +
0.02 - 0.04) = 80.4 / 0.06 \approx 1{,}340$ model units
($\pounds 134$M), close to the observed $\pounds 138.5$M.  The
first-order prediction was $\pounds 200$M; the secondary sinks
account for the gap.

The qualitative implication is unchanged.  The wealth tax must
exceed $r$ to produce a finite analytical cap; the secondary sinks
contribute approximately $\delta_d + \iota/L \approx 0.026$ of
effective tax, so a nominal wealth tax of $\tau \approx 0.014$
would suffice to make the cap analytically finite (rather than
$\tau \approx 0.04$ from the first-order argument).  All real-world
proposals modelled have nominal $\tau \geq 0.01$ at their lowest
tier; under the tightened analysis several of them therefore enter
the analytically-capped regime at sufficiently high thresholds.
The reading of which policies rely on cap-by-mortality and which
do not changes correspondingly.

\subsection{The cap-by-mortality regime}

At the default capital return $r = 0.04$, Zucman's 2\% above
\pounds 10M places the lowest active bracket at $\tau = 0.02 < r$.
The analytical fixed-point is therefore not finite, and the
empirical $\approx \pounds 270$M cap on the top wealth observed in
simulation is set by demographic churn rather than by tax.  We refer
to this as the \emph{cap-by-mortality regime}: a wealth-tax design
whose lowest active rate is below the rate of capital return relies
on mortality and secondary sinks for the analytical cap, which is a
fragile basis for a policy claim.

Schemes whose lowest tier rate exceeds $r$ are structurally robust
because they cap the top analytically.  By that test, none of the
real-world wealth-tax proposals we model (Zucman 2\%, Warren 2\% at
the lowest bracket, Sanders 1\%, Piketty 0.1\%) has $\tau > r$ at
their lowest bracket under default $r = 0.04$.  At the calibrated
$r = 0.06$ (Section~\ref{sec:calibration}) the gap widens and the
cap-by-mortality regime becomes more salient.

\section{Calibration to UK wealth distribution}
\label{sec:calibration}

The wealth-tax-design conclusions in this paper depend on the
assumed value of $r$.  We attempt a partial calibration against UK
data.

\subsection{The calibration target}

UK Sunday Times Rich List 2024 reports approximately 165 sterling
billionaires; the wealth of the richest is approximately
\pounds 37B.  A Pareto fit to the rich-list distribution gives slope
$-0.706$ on $\log(\text{rank})$ vs $\log(\text{wealth})$.  WID
estimates of top wealth shares \citep{saez2016wealth} give
top-1\% share $0.30$, top-10\% share $0.57$, top-0.1\% share
$0.13$.  We use these four numbers as calibration targets.

\subsection{The grid search}

We sweep over capital return $r \in \{0.04, 0.06, 0.08, 0.10\}$,
rent share $\in \{0.20, 0.35, 0.50\}$, and inheritance tax
$\iota \in \{0.0, 0.2, 0.4\}$, holding all other dials at default
plus a Tory-baseline spending package (no UBI, no state employment).
180 cells, 5 seeds at $N = 1000$, 200 ticks.  The loss function is
the L2 distance to the four targets, with the slope target weighted
$0.1$ relative to the share targets.

The best L2 fit is at $r = 0.06$, rent share $0.50$, inheritance
tax $0.40$, yielding observed top-1\% $0.15$, top-10\% $0.65$,
top-0.1\% $0.06$, Pareto slope $-0.673$.  The slope matches the UK
fit within $5\%$; the absolute levels are too low by approximately
a factor of two at top-1\% and top-0.1\%.

\subsection{Joint calibration with subsistence and bequest cap}
\label{sec:joint-cal}

The grid search above held the subsistence cost $\sigma$ and the
inheritance cap $\bar{w}_\text{inh}$ at default ($0.6$ and
$\infty$ respectively) and matched the wealth-tail targets at the
cost of catastrophic mortality.  A wider grid that includes
$\sigma \in \{0.3, 0.6\}$ and $\bar{w}_\text{inh} \in \{50, 500,
\infty\}$ recovers a substantially better joint fit.

\begin{table}[h]
\centering
\small
\caption{Top 10 best-fit cells in the joint calibration ranking,
with subsistence $\sigma$ and inheritance cap $\bar{w}_\text{inh}$
added to the search space.  Total grid: 108 configurations, 4 seeds
each.  Loss is the weighted L2 against the five UK targets including
the alive-rate target.}
\label{tab:joint-cal}
\begin{tabular}{rrrrrrrrrrr}
\toprule
rank & $r$ & rent & $\iota$ & $\sigma$ & $\bar{w}_\text{inh}$
& top-1 & top-10 & top-0.1 & slope & alive \\
\midrule
  1 & 0.08 & 0.20 & 0.40 & 0.30 &  50 & 0.21 & 0.58 & 0.07 & $-0.631$ & 0.80 \\
  2 & 0.08 & 0.35 & 0.00 & 0.30 &  50 & 0.21 & 0.60 & 0.07 & $-0.635$ & 0.78 \\
  3 & 0.08 & 0.20 & 0.00 & 0.60 &  50 & 0.23 & 0.62 & 0.07 & $-0.648$ & 0.72 \\
  4 & 0.08 & 0.50 & 0.00 & 0.30 &  50 & 0.24 & 0.67 & 0.07 & $-0.719$ & 0.74 \\
  5 & 0.08 & 0.35 & 0.40 & 0.30 &  50 & 0.25 & 0.70 & 0.09 & $-0.631$ & 0.77 \\
  6 & 0.08 & 0.20 & 0.00 & 0.30 &  50 & 0.16 & 0.52 & 0.04 & $-0.602$ & 0.82 \\
  7 & 0.06 & 0.35 & 0.40 & 0.30 &  50 & 0.15 & 0.56 & 0.04 & $-0.509$ & 0.77 \\
  8 & 0.06 & 0.20 & 0.00 & 0.30 & 500 & 0.17 & 0.68 & 0.03 & $-0.475$ & 0.82 \\
  9 & 0.06 & 0.50 & 0.00 & 0.30 &  50 & 0.16 & 0.54 & 0.04 & $-0.582$ & 0.74 \\
 10 & 0.08 & 0.20 & 0.40 & 0.60 &  50 & 0.22 & 0.68 & 0.07 & $-0.641$ & 0.71 \\
\bottomrule
\end{tabular}
\end{table}

\paragraph{The headline.}  The best joint fit (rank 1) has $r =
0.08$, rent $= 0.20$, $\iota = 0.40$, $\sigma = 0.30$,
$\bar{w}_\text{inh} = 50$.  This yields top-10 share 0.58 (matching
the UK target of 0.57 within 0.01), Pareto slope $-0.631$ (within
$10\%$ of $-0.706$), top-1 share 0.21 (below the UK 0.30), top-0.1
share 0.07 (below the UK 0.13), and \emph{alive rate 0.80}.

The 80\% alive rate is a substantial improvement over the
26\% reported under the narrower calibration of
\S\ref{sec:calibration}.  The configuration is therefore
demographically plausible.

\paragraph{Mechanisms.}  Two of the three new dials do the work.

\emph{Lower subsistence $\sigma = 0.3$} reduces the mortality cost
of high extraction by giving each agent more headroom against the
deterministic per-tick subsistence drain.  This is a parameter that
arguably should not be free in a calibration (subsistence is an
empirical quantity), but for the model's stylised economy the
appropriate value is uncertain.

\emph{Lower inheritance cap $\bar{w}_\text{inh} = 50$
model units (\pounds 5M)} truncates dynastic accumulation directly
without requiring a high inheritance tax to do so.  The cap acts
as an absolute upper bound on what a single agent can inherit at
death.  This finding is novel relative to the existing wealth-tax-
ABM literature: the dial that matters most for matching the
UK tail while preserving survival is the bequest cap, not the
inheritance tax \emph{rate}.

\paragraph{Implications for the policy comparison.}  The Pareto
slope is reachable at a configuration with realistic demographics.
The remaining tail-share mismatch (top-1 and top-0.1 are too low)
suggests the next calibration step should free the population
proportions (the rentier share at 3\% of the population is
plausibly too high for the strict billionaire class).  A full
joint calibration with population proportions free is the natural
follow-up paper and would let policy claims be expressed in pounds
without caveat.

\subsection{The mortality--tail trade-off}

The calibrated configuration keeps only 26\% of the population alive
at terminal tick.  At UK life expectancy and current population
size, the observed survival rate over a 250-year window would be
substantially higher.  The model cannot currently match Pareto slope
and demographic survival simultaneously; raising $r$ and rent to
match the wealth tail accelerates extraction past the rate at which
the population can replace itself.  This is the central remaining
calibration limitation of the model.

We do not claim the calibrated configuration as a fully validated
representation of UK reality.  We claim that the Pareto slope is
\emph{reachable}, which is enough to use the calibrated parameters
as a sensitivity case for the policy comparisons.

%==========================================================================
\section{Policy results: default vs calibrated, side by side}
\label{sec:policies}
%==========================================================================

We rerun the policy comparison at $N = 10{,}000$ under both default
and calibrated parameters, holding the policy-distinguishing dials
(tax schedule and policy-specific spending package) fixed across
policies.  Tables~\ref{tab:policies-default} and
\ref{tab:policies-calibrated} report the comparison side by side.

\begin{table}[t]
\centering
\small
\caption{Eight wealth-tax policies at default $r = 0.04$,
rent share $0.35$.  $N = 10{,}000$, 10 seeds, 180 ticks.  State
excluded from top-1 share.}
\label{tab:policies-default}
\begin{tabular}{lrrrrrr}
\toprule
policy & alive & $\Phi$ & top-1 share & $\geq$\pounds10M & $\geq$\pounds100M & $\geq$\pounds1B \\
\midrule
uk\_tory\_2010\_24            & 6{,}268 & 0.118 & 0.183 & 1{,}540 & 796 & \textbf{19} \\
uk\_labour\_2024              & 6{,}235 & 0.288 & 0.167 & 1{,}435 & 453 & 1 \\
zucman\_2pct\_hoard           & 6{,}103 & 0.438 & 0.106 & 1{,}440 & 27  & 1 \\
zucman\_2pct\_spend           & 8{,}053 & 0.440 & 0.095 & 1{,}762 & 69  & 1 \\
zucman\_progressive\_2\_4\_8 & 8{,}059 & 0.446 & 0.087 & 1{,}761 & 64  & 1 \\
warren\_2020                  & 7{,}699 & 0.310 & 0.103 & 1{,}749 & 704 & 1 \\
piketty\_book                 & 7{,}777 & 0.432 & 0.135 & 1{,}456 & 108 & 1 \\
reform\_uk\_flavour           & 5{,}334 & 0.060 & 0.172 & 1{,}632 & 768 & \textbf{97} \\
\bottomrule
\end{tabular}
\end{table}

\begin{table}[t]
\centering
\small
\caption{Same eight policies at calibrated $r = 0.06$, rent share
$0.50$.  $N = 10{,}000$, 6 seeds, 180 ticks.  State excluded.}
\label{tab:policies-calibrated}
\begin{tabular}{lrrrrrr}
\toprule
policy & alive & top-1 share & $\geq$\pounds10M & $\geq$\pounds100M & $\geq$\pounds1B & max wealth \\
\midrule
uk\_tory\_2010\_24            & 2{,}230 & 0.198 & 866 & 309 & \textbf{300} & \pounds48B \\
uk\_labour\_2024              & 3{,}029 & 0.157 & 1{,}360 & 800 & 262 & \pounds12B \\
zucman\_2pct\_hoard           & 2{,}532 & 0.262 & 910 & 306 & 10 & \pounds4.1B \\
zucman\_2pct\_spend           & 7{,}734 & 0.233 & 1{,}746 & 801 & 46 & \pounds2.4B \\
zucman\_progressive\_2\_4\_8  & 7{,}744 & 0.143 & 1{,}750 & 720 & \textbf{0} & \pounds0.4B \\
warren\_2020                  & 7{,}234 & 0.179 & 1{,}740 & 804 & 140 & \pounds1.5B \\
piketty\_book                 & 7{,}363 & 0.230 & 1{,}441 & 800 & 20 & \pounds1.7B \\
reform\_uk\_flavour           & 2{,}526 & 0.181 & 1{,}010 & 341 & \textbf{324} & \pounds77B \\
\bottomrule
\end{tabular}
\end{table}

Three things move materially.

\paragraph{Billionaire counts grow under most policies.}  Tory now
produces 300 billionaires (vs 19 at default), Reform 324 (vs 97).
Even Zucman 2\% flat with spending leaves 46 billionaires (vs 1).
The model at calibrated parameters is much closer to the UK Rich
List distribution (Tory max wealth \pounds 48B vs UK Rich List top
of approximately \pounds 37B) but the per-capita billionaire density
now \emph{over}-shoots in the other direction.

\paragraph{Mortality is much higher in non-spend policies.}  Tory
keeps 22\% of the population alive at $r = 0.06$, vs 63\% at
$r = 0.04$.  The combination of higher $r$ and higher rent share
extracts faster than the population can replace.  This is the
calibration--mortality trade-off made concrete across policies.

\paragraph{Zucman progressive binds and Zucman flat does not.}
Under calibrated parameters Zucman progressive produces \emph{zero}
billionaires while Zucman flat leaves 46, a clean separation that
did not exist at default $r = 0.04$.  Top-1 share also separates
clearly: $0.143$ progressive vs $0.233$ flat.  The earlier ``top
tier does not bind'' reading was specific to the default-parameter
regime where wealth does not reach the \pounds 1B bracket; once
$r$ is raised to the calibration target, the upper tier becomes the
decisive instrument.

\subsection{Implications for the fixed-point reading}

At calibrated $r = 0.06$, none of the eight policy lowest-tier rates
exceed $r$ (Zucman 2\%, Warren 2\%, Sanders 1\%, Piketty 0.1\%).
The cap-by-mortality regime identified in
Section~\ref{sec:fixedpoint} now bites harder.  The Zucman
progressive top tier of 4\% above \pounds 100M (stacking to 8\%
above \pounds 1B) is the one tier among the modelled policies that
exceeds the calibrated $r$, and it is the one policy that eliminates
billionaires.  The \pounds 1B tier is no longer a tautological
non-binding constraint; it is the mechanism by which the policy
succeeds in this regime.  The $\tau$-vs-$r$ design test is
preserved and becomes more rather than less significant under
calibration.

%==========================================================================
\section{Cap and tax as complementary instruments: the ownership-
redistribution result}
\label{sec:ownership}
%==========================================================================

The fixed-point analysis in Section~\ref{sec:fixedpoint} and the
policy-comparison Tables~\ref{tab:policies-default} and
\ref{tab:policies-calibrated} characterise the conditions under
which the upper tier of a progressive wealth tax binds.  The
broader question of \emph{what each fiscal instrument does for
the wealth distribution as a whole} requires a different
comparison.  We report it here.

At the joint-calibrated UK baseline, we compare five interventions:
status quo, Zucman 2\% wealth tax above \pounds 10M, a ban on
interest on loans, a \pounds 5M absolute cap on inheritance, and
the cap combined with Zucman 2\%.  The principal outcome variables
are the wealth share held by the bottom 50\% of the population
(the \emph{ownership stake}) and the cumulative state revenue
raised across the run.  Table~\ref{tab:cap-tax-complementary}
reports the comparison.

\begin{table}[t]
\centering
\small
\caption{Five interventions at the joint-calibrated UK baseline.
$N = 10{,}000$, 4 seeds, 180 ticks.  Cap and tax are
complementary instruments: the tax raises revenue, the cap
redistributes ownership.}
\label{tab:cap-tax-complementary}
\begin{tabular}{lrrrr}
\toprule
intervention & bottom-50\% share & top-10\% share
& cumulative state revenue & annual revenue \\
\midrule
status quo                & 0.0\%  & 100.0\% & \pounds 0      & \pounds 0     \\
Zucman 2\% above \pounds 10M & 0.6\%  &  98.9\% & \pounds 5{,}165B & \pounds 56B    \\
interest ban ($i = 0$)    & 0.0\%  & 100.0\% & \pounds 0      & \pounds 0     \\
\pounds 5M bequest cap     & \textbf{19.2\%} &  54.0\% & \pounds 484B  & indirect    \\
cap + Zucman 2\%           & \textbf{25.3\%} &  39.9\% & \pounds 335B  & \pounds 0.8B  \\
\bottomrule
\end{tabular}
\end{table}

The result has a clean economic interpretation.

\paragraph{The wealth tax is a revenue instrument.}  Zucman's 2\%
above \pounds 10M raises approximately \pounds 56B per year at
simulated UK scale and accumulates approximately \pounds 5{,}165B
in state wealth over the 180-tick run.  The state can then
recirculate this revenue via the spending package
(UBI, state employment, state purchases) and the recirculation
improves survival and welfare flows for the bottom half.  But the
recirculation is a flow; it does not give the bottom half an
ownership stake in the underlying capital.  The bottom-50\% share
rises from 0.0\% to 0.6\% under Zucman alone --- essentially noise.

\paragraph{The cap is an ownership-redistribution instrument.}
The \pounds 5M absolute cap raises far less revenue through the
explicit wealth-tax channel (zero, when paired with no wealth tax)
but accumulates \pounds 484B of state wealth via inheritance
excess at death.  More importantly, the cap raises the bottom-50\%
wealth share from 0.0\% to 19.2\%, and to 25.3\% when combined
with the wealth tax.  The mechanism is direct: dynastic wealth
above the cap is either redirected to the state (and then
distributed via UBI) or simply does not propagate to the next
generation.  In either case the stock of dynastically-held wealth
shrinks and the bottom half can hold a meaningful share for the
first time.

\paragraph{The two instruments compete for the upper-end tax
base.}  Cap-plus-Zucman raises less wealth-tax revenue than Zucman
alone (\pounds 0.8B/year vs \pounds 56B/year) because the cap
prevents the buildup of large wealth above the wealth-tax
threshold that would otherwise have been taxed.  This is the
single trade-off the joint use of the two instruments produces.
The trade-off is real but small relative to the structural
benefit of redistributed ownership.  An economy with cap-plus-tax
has 25\% of its wealth held by the bottom half of the population,
versus 0.6\% under tax alone; the foregone tax revenue (\pounds 55B
per year in the simulation) is the price paid for that
redistribution.

\paragraph{The interest ban is neither.}  Banning interest on
loans does not raise revenue and does not redistribute ownership.
The bottom-50\% share is unchanged.  The intervention improves
working-class survival slightly (88\% vs 84\% under status quo)
because workers do not compound debt under zero interest, but the
distributional impact at the top is zero.  The historical and
religious case for interest bans is largely an ethical case about
debt and dependence on the borrower side; it is not, by this
model's read, an instrument for moving the wealth distribution.

\paragraph{Choice of instrument depends on the policy goal.}  If
the goal is funding the state at scale, use the wealth tax.  If
the goal is redistributing the ownership of capital, use the cap.
If the goal is both, use both, and accept the trade-off in
wealth-tax revenue.  The contemporary debate has foregrounded the
revenue instrument; this paper foregrounds the ownership
instrument, in part because the ownership instrument is the one
that is absent from the contemporary debate.

%==========================================================================
\section{Comparison to Spangler \& Sarkar}
\label{sec:spangler}
%==========================================================================

\citet{spangler2022distribution} is the closest published ABM
neighbour and finds that skill inheritance raises wealth inequality.
We can test their mechanism directly in our model under a clean
toggle that holds bandwidth-gating and class-bias persistence intact
and swaps only whether transmitted skill values are parent-derived
or class-typical.  A coarser toggle that simultaneously destroys
bandwidth and class-bias persistence (the mutation rate $\mu = 1$
arm) is reported for contrast.

\begin{table}[h]
\centering
\small
\caption{Spangler \& Sarkar test at $N = 100$, 10 seeds, 250 ticks.}
\label{tab:spangler}
\begin{tabular}{lrrr}
\toprule
metric & default & class-prior & no-inheritance \\
       & (parent-derived) & (clean toggle) & ($\mu = 1$, confounded) \\
\midrule
Gini             & 0.885 & 0.876 & 0.908 \\
Top-1 share      & 0.235 & 0.250 & 0.254 \\
Top-10 share     & 0.785 & 0.683 & 0.891 \\
Alive            & 61.5  & 56.5  & 73.0 \\
\bottomrule
\end{tabular}
\end{table}

Under the clean toggle, parent-derived skill values raise Gini by
$0.010$ relative to a class-typical prior.  The direction is a
partial reproduction of \citet{spangler2022distribution}; the
magnitude is much smaller than they report.  The mechanism for the
small magnitude is class-specific skill-to-wealth elasticity
(documented in the ALIFE companion submission): skill inheritance
acts strongly on workers and makers, but the wealth dynamics that
dominate the overall Gini are anchored in rentier and lender stocks
where skills have weak elasticity.  The inheritance-to-inequality
link is mechanism-specific and not robust to the choice of ABM
substrate.

%==========================================================================
\section{Discussion}
%==========================================================================

\subsection{What we believe is novel for a JEIC audience}

The analytical $w^* = \tau \theta / (\tau - r)$ fixed-point
characterisation (Section~\ref{sec:fixedpoint}) is, to our reading
of the wealth-tax-ABM literature, not previously reported.  It is a
first-order rigorous statement that puts a sharp condition on
the relationship between the rate of capital return and the lowest
active tax bracket.  \citet{benhabib2011distribution} provide the
heterogeneous-agent-macro neighbour and our result recovers the
qualitative structure of theirs in an ABM substrate.

The calibration-sensitivity result
(Section~\ref{sec:policies}) is the second contribution.  The
``top tier does not bind'' conclusion at default parameters is
reversed at the calibrated parameters; the Zucman progressive
schedule is mechanically idle at one parameter setting and the
load-bearing instrument at the other.  The policy conclusion of the
wealth-tax-design analysis is therefore parameter-conditional, and
we report it as such.

The Spangler \& Sarkar reproduction
(Section~\ref{sec:spangler}) is the third.  The clean toggle gives a
partial reproduction in direction with a much smaller magnitude than
the single-channel ABM produces.  The inheritance-to-inequality
link is mechanism-specific.

\subsection{What the model cannot do}

The Pareto-slope calibration is partial: matching the slope requires
parameters that produce 26\% survival, well below UK reality.
A joint calibration that frees the population proportions and
demographic parameters to fit both the wealth tail and the survival
rate is the natural follow-up paper.

The model has no monetary-policy layer, no debt cycle in the formal
sense, no firm formation, no spatial structure, and a single lender
pool with no shadow-banking.  The Spangler test is run in our
multi-channel setting and may not reproduce in their single-channel
design; the qualitative-direction reproduction we report is the only
empirical contact we have with their model.

\subsection{Honest reading}

The wealth-tax-design conclusion of this paper is:
\emph{the upper-tier rates of a progressive wealth-tax schedule are
mechanically idle or load-bearing depending on the assumed rate of
capital return}.  At UK default-parameter values they are idle.  At
parameters consistent with the UK Pareto slope they are load-bearing.
A policy designer should choose the rate of capital return their
analysis is conditional on and report the conclusion as conditional;
we have done both.  We do not adjudicate which $r$ is the right
choice for the UK because that question is empirical and our
calibration only partially settles it.

%==========================================================================
\section{Conclusion}
%==========================================================================

We have derived an analytical fixed-point characterisation of top
wealth under a flat marginal wealth tax,
$w^* = \tau \theta / (\tau - r)$ for $\tau > r$.  We have validated
it against an agent-based simulation and shown the
cap-by-mortality regime that obtains for $\tau \leq r$.  We have
partially calibrated the model to UK wealth-distribution data and
shown the calibration-mortality trade-off.  We have rerun the
policy-design comparison at default and calibrated parameters and
shown the calibration sensitivity of the central finding.  And we
have engaged \citet{spangler2022distribution} via a clean-toggle
reproduction and reported a partial reproduction in direction with
a much smaller magnitude than their single-channel ABM.

The policy implication, presented as conditional, is: under
default-parameter conditions the lower bracket of a progressive
wealth tax sets the cap and the upper tiers are mechanically idle;
under conditions consistent with the empirical UK wealth tail, the
upper tiers are load-bearing and the Zucman progressive schedule
sharply outperforms the flat 2\% form.

\bibliographystyle{plainnat}
\bibliography{references}

\appendix
\section{Dial reference}
\label{app:dials}

Default values for the dial panel are listed below.  Tax thresholds
are in model units; one unit corresponds to \pounds 100{,}000 for
this paper's policy-comparison content.

\begin{table}[h]
\centering
\small
\begin{tabular}{lrl}
\toprule
dial & default & group \\
\midrule
wage                          & 1.0 & production \\
maker\_multiplier             & 1.6 & production \\
employer\_markup              & 0.4 & production \\
capital\_return ($r$)         & 0.04 & capital \\
rent\_share\_of\_wage         & 0.35 & capital \\
interest\_rate ($i$)          & 0.06 & capital \\
wealth\_tax\_tiers            & ()   & state \\
ubi                           & 0.0  & state \\
state\_employs\_fraction      & 0.0  & state \\
state\_buys\_maker\_output    & 0.0  & state \\
subsistence                   & 0.6  & life \\
lifespan\_mean / std          & 70 / 12 & life \\
inheritance\_tax\_rate        & 0.0  & transmission \\
population\_scale             & 1    & logistics \\
seed                          & 42   & logistics \\
\bottomrule
\end{tabular}
\end{table}

\end{document}
